%%%%%%%%%%%%%%%%%%%%%%%%%%%%%%%%%%%%%%%%%%%%%%%%%%%%%%%%%%%%%%%%%%%%%%%%%%%%
%% Trim Size: 9.75in x 6.5in
%% Text Area: 8in (include Runningheads) x 5in
%% ws-ijseke.tex     22-2-07
%% Tex file to use with ws-ijseke.cls written in Latex2E.
%% The content, structure, format and layout of this style file is the
%% property of World Scientific Publishing Co. Pte. Ltd.
%% Copyright 1995, 2002 by World Scientific Publishing Co.
%% All rights are reserved.
%%%%%%%%%%%%%%%%%%%%%%%%%%%%%%%%%%%%%%%%%%%%%%%%%%%%%%%%%%%%%%%%%%%%%%%%%%%%
\documentclass[final]{ws-ijseke}
\usepackage{ctex}

\begin{document}


%%%%%%%%%%%%%%%%%%%%% Publisher's Area please ignore %%%%%%%%%%%%%%%
%
%\catchline{}{}{}{}{}
%
%%%%%%%%%%%%%%%%%%%%%%%%%%%%%%%%%%%%%%%%%%%%%%%%%%%%%%%%%%%%%%%%%%%%
\pagestyle{empty}
\title{$ELBO$（证据下界）综述 }

\author{聂宇航}

\address{信息科学与工程学部, 中国海洋大学, 山东 青岛\\
\email{nieyuhang@stu.ouc.edu.cn}
}


\maketitle


\begin{abstract}
摘  要：证据下界（Evidence Lower Bound, ELBO）是变分推断（Variational Inference）$^{[1]}$的理论基石，通过优化可处理的变分分布逼近复杂后验分布，解决概率模型中隐变量推断的难解性问题。本文系统综述$ELBO$的数学本质、推导方法、优化策略与应用前沿。
\end{abstract}

\keywords{证据下界（ELBO）; 变分推断; 最大似然估计（MLE）; KL 散度; 重参数化技巧; 隐变量模型; 变分自编码器（VAE）; 贝叶斯神经网络（BNN）; 后验近似; 蒙特卡洛估计; 重要性加权自编码器（IWAE）;}


\section{ 引言 \label{ga}}	%) A SECTION HEADING

\hspace*{2em} 在当代机器学习与统计建模中，复杂概率模型的构建已成为刻画高维不确定性数据的重要手段。然而，在这类模型中往往涉及隐变量（latent variables）$^{[2]}$，如图模型中的隐状态、生成模型中的潜在表示等。这些隐变量使得后验分布 $p(z|x)$ 难以精确计算，尤其在缺乏解析解或需要高维积分的情形下，后验推断成为学习中的主要瓶颈。\\
\hspace*{2em} 为了解决这一难题，研究者提出了多种近似推断方法。其中，变分推断（Variational Inference, VI）作为一种高效的优化方法，近年来受到广泛关注。它将后验推断问题转化为变分分布 $q(z)$与真实后验 $p(z|x)$ 之间的距离最小化问题。具体而言，VI 通过引入可调变分族Q，用优化替代采样，避免了马尔可夫链蒙特卡洛（MCMC）方法常见的收敛性难题与采样效率低下的问题。\\
在变分推断中，证据下界（Evidence Lower Bound, $ELBO$）作为核心目标函数，承载着连接观测数据似然 $p(x)$、模型联合分布$p(x,z)$、以及变分分布 $q(z)$ 的重要作用。$ELBO$ 的最大化不仅是推断质量的度量标准，更是近似后验与生成模型的学习信号。最大化 $ELBO$ 等价于最小化$KL(q(z)||p(z|x))$，从而实现模型学习与推断质量的统一。\\
\hspace*{2em} 随着深度学习的兴起，$ELBO$ 被广泛应用于变分自编码器（VAE）、贝叶斯神经网络、主题模型、图神经网络$^{[3]}$等领域，成为现代生成建模与不确定性量化的理论基石。此外，围绕 $ELBO$ 的优化方法与理论改进也不断发展，如重参数化技巧、重要性加权下界（IWAE）、$\chi$ 散度下界、半摊销推断（semi-amortized inference）$^{[4]}$等，极大地扩展了其适用范围和表达能力。\\
\hspace*{2em} 本文旨在系统梳理 $ELBO$ 的理论基础、数学推导、物理意义与优化技术，并总结其在深度学习与贝叶斯建模中的典型应用与研究进展。通过对 $ELBO$ 的全面综述，我们希望为读者提供一个统一的理解框架，帮助其把握这一核心工具的本质与发展趋势。

\section{$ELBO$ 的数学基础与推导$^{[5-7]}$ \label{sec1}}
\subsection{问题定义与建模背景}

\hspace*{2em}我们最主要的目标是对一个含隐变量的概率模型进行参数估计。设可观察变量为 $x$，隐变量为 $z$，待估计的模型参数为 $\theta$。在经典的情形下，我们常使用最大似然估计（MLE）对 $\theta$ 进行估计，但由于存在无法观测的隐变量 $z$，直接计算似然函数变得困难。为此，EM算法应运而生，其本质仍然是最大似然，只是在每一次 M 步之前，先用 E 步估计并“固定”住隐变量的分布。
现在我们要对观察变量做最大似然估计，即最大化以下函数：
\begin{equation*}
L(\theta) = \int_{x\sim p(x)}{logp_\theta(x)}
\end{equation*}

由于 $x$ 与 $z$ 联合生成，因此 $p_\theta(x)$ 需要对所有可能的隐变量 $z$ 积分（或求和）边缘化：
\begin{equation*}
logp_\theta(x) = log\int p_\theta(x,z)dz
\end{equation*}
\hspace*{2em} 但由于模型中有隐变量，这一积分在很多模型中并不能解析计算，所以我们必须把隐变量确定下来，这样就可以正常地做MLE了。至于如何根据观察把隐变量确定下来，我们的思路也很容易想到，即使用贝叶斯推断：
\begin{equation*}
p(z|x) = p(z) \frac{p(x|z)}{p(x)}
\end{equation*}
\hspace*{2em} 得到隐变量的后验分布$p(z|x)$之后，我们并不打算去估计出某个确定的值，因为我们认为这样会损失更多信息。我们的打算是把整个分布都利用起来，即：
\begin{equation*}
L(\theta)= \mathbb{E}_{x\sim p_{data}(x)}[logp_\theta(x)] = \mathbb{E}_x[L_x(\theta)] = \mathbb{E}_x[logp_\theta(x)]
\end{equation*}
其中:
\begin{align*}
L_x(\theta) &= logp_\theta(x) \\ &= log(\mathbb{E}_z[p_\theta(x,z)]) \\ &=   log(\mathbb{E}_z[p_\theta(z|x)p_\theta(x)]) \\&= log(\mathbb{E}_z[p_\theta(z|x)\frac{p_\theta(x, z)}{p_\theta(z|x)}])
\\ &= log\mathbb{E}_{z\sim q(z|x)}(\frac{p_\theta(x,z)}{q(z|x)})
\end{align*}
\hspace*{2em} 在上式的第二个等号，我使用乘法公式（贝叶斯公式）分解出$z$ 的后验分布；而在第三个等号，我再次使用乘法公式把$p_\theta(x)$ 分解为$\frac{p_\theta(x,z)}{p_\theta(z|x)}$来计算，因为$p_\theta(z|x)$没法儿算。看上去，这就好像在$p_\theta(x, z)$上做了一个同时乘和除$p_\theta(z|x)$的数学技巧，但它实际上是有动机的。把一个难计算的 $p_\theta(x)$ 变成了一个期望形式，这个期望是相对于我们自己定义的分布 $q(z|x)$，是可采样和可学习的。\\
\hspace*{2em} 但$p_\theta(z|x)$ 是难以计算的。
例如，在VAE中，这儿的$\theta$代表的是解码器的参数。根据解码器输出高斯分布的假设，显然我们可以显式地计算出$p_\theta(x|z)$ ，而我们又有解码器输入、即隐变量先验为标准高斯分布的假设，所以也可以显式地写出$p_\theta(z)$ （就是标准高斯），然而我们无法从解码器中得到$p_\theta(z|x)$  。\\\hspace*{2em} 在数学上看$p(z|x) = \frac{p(x,z)}{p(x)}$,其中$p(x) = \int p(x,z)dz = \int p(x|z)p(z)dz$ 通常是一个高阶积分，难以求解。
\\ \hspace*{2em} 因此为了另辟蹊径，我们使用变分推断来近似计算$p_\theta(z|x)$ 。
 \subsection{从 KL 散度视角导出 $ELBO$ $^{[8-10]}$}
 \hspace*{2em} 对于我们的目标$p_\theta(z|x)$ ，我们要找一个最好的$q_\phi(z|x)$来近似它。怎么样的近似才好呢，标准就是最小化
 $KL$散度，即$KL(q_\phi(z|x)||p_\theta(z|x))$\\
 \hspace*{2em} 那为什么是KL散度呢$KL(q(z)||p(z)) = \int q(z)log\frac{q(z)}{p(z)}dz$,而不是我们可以轻松的想到$\int (q(z)−p(z))^2dz$或者$\int |q(z)−p(z)|dz$呢?\\ \hspace*{2em} 因为$KL$散度代表：如果我们用 $q(z)$ 来编码数据，而数据实际上来自 $p(z)$，我们每单位多用了多少“信息”。这直接连到压缩理论、通信理论，和“最优表示”的理论目标一致。并且KL 散度对常见分布（如高斯）有解析解，以及它是光滑的、可导的，能反向传播优化。
 如果把$q_\phi(z|x))$看成自变量函数，那么这个散度就是关于它的一个泛函，于是我们现在正是在求一个泛函极值问题。\\
 现在来求这个泛函极值问题：\\
\begin{equation*}
KL(q_\phi(z|x)||p_\theta(z|x)) = \int_z q_\phi(z|x)log\frac{q_\phi(z|x)}{p_\theta(z|x)}dz
\end{equation*}
\hspace*{2em} 问题是我们的目标就是 $p_\theta(z|x)$，积分里还有它、肯定没法求。得把它拆开，还是用乘法公式，就得到：
 \begin{equation*}
\int_z q_\phi(z|x)log\frac{q_\phi(z|x)p_\theta(x)}{p_\theta(x,z)}dz
\end{equation*}
再化个简，把跟积分变量$z$无关的项$p_\theta(x)$剥离出来:
 \begin{align*}
KL(q_\phi(z|x)||p_\theta(z|x)) &=\int_z q_\phi(z|x)log\frac{q_\phi(z|x)p_\theta(x)}{p_\theta(x,z)}dz\\ & = \int_z q_\phi(z|x)log\frac{q_\phi(z|x)}{p_\theta(x,z)}dz + \int_z q_\phi(z|x)logp_\theta(x)dz \\ &= \int_z q_\phi(z|x)log\frac{q_\phi(z|x)}{p_\theta(x,z)}dz + logp_\theta(x)\int_z q_\phi(z|x)dz\\ &= \int_z q_\phi(z|x)log\frac{q_\phi(z|x)}{p_\theta(x,z)}dz + logp_\theta(x)
 \end{align*}
 改写为以似然为中心：
 \begin{equation*}
L_x(\theta) = logp_\theta(x) = \int_z q_\phi(z|x)log\frac{p_\theta(x,z)}{q_\phi(z|x)}dz + KL(q_\phi(z|x)||p_\theta(z|x))
\end{equation*}
左边这项$\int_z q_\phi(z|x)log\frac{p_\theta(x,z)}{q_\phi(z|x)}dz = \mathbb{E}_{q_\phi (z|x)}[log\frac{p_\theta(x,z)}{q_\phi(z|x)}]$就是$ELBO$。让我们继续看看它的内涵：
 \begin{align*}
ELBO &= \int_zq_\phi(z|x)log\frac{p_\theta(x,z)}{q_\phi(z|x)}dz \\ &= \int_zq_\phi(z|x)log\frac{p_\theta(x|z)p_\theta(z)}{q_\phi(z|x)}dz \\ &=
\mathbb{E}_{q_\phi (z|x)}[logp_\theta(x|z)] + \mathbb{E}_{q_\phi (z|x)}[log\frac{p_\theta(z)}{q_\phi(z|x)}] \\ &=
\mathbb{E}_{q_\phi (z|x)}[logp_\theta(x|z)] - KL(q_\phi(z|x)||p_\theta(x))
\end{align*}
\hspace*{2em} 现在， $L_x(\theta)$中除了$KL(q_\phi(z|x)||p_\theta(z|x))$项，其他地方都不再包含$p_\theta(z|x)$ ，取而代之的是编码器
$q_\phi(z|x)$。所以，我们干脆丢掉$KL(q_\phi(z|x)||p_\theta(z|x))$，只最大化$ELBO$，从而不再有任何计算障碍。事实上，由于
 $KL(q_\phi(z|x)||p_\theta(z|x)) \ge 0$（等号成立当且仅当$p_\theta(z|x) = q_\phi(z|x)$），所以我们有：
 \begin{equation*}
L_x(\theta) = logp_\theta(x) \ge ELBO
\end{equation*}
\hspace*{2em} 所以，我们实际上是在最大化似然的下界。实际上，因为$p_\theta(x)$又叫证据，所以$ELBO$叫做$ELBO$（证据下界）。
然而，直接丢弃 $KL(q_\phi(z|x)||p_\theta(z|x))$是合理的吗？是合理的，而合理性本质上是由EM算法保证的。
\subsection{从 Jensen 不等式推导 $ELBO$ $^{[11-12]}$}
$ELBO$ 也可以通过 Jensen 不等式直接从$logp_\theta(x)$的定义导出：
\begin{align*}
    L_x(\theta) = logp_\theta(x) &= log\int_z p_\theta(x,z)dz \\ &= log\int_z \frac{q_\phi(z|x)}{q_\phi(z|x)}p_\theta(x,z)dz \,\,\, (p_\theta(x,z) = p_\theta(x|z)p(z))\\&= log \int_z q_\phi(z|x)\frac{p_\theta(x|z)p(z)}{q_\phi(z|x)}dz\,\,\, (\text{对数是凹函数，使用Jensen不等式})\\&\ge
    \int_z q_\phi(z|x)log\frac{p_\theta(x|z)p(z)}{q_\phi(z|x)}dz \\&=  \mathbb{E}_{x\sim q_\phi(z|x)}[\frac{logp_\theta(x,z)}{q_\phi(z|x)}]  =ELBO = \mathcal L\\&=
    \int_z q_\phi(z|x)logp_\theta(x|z)dz + \int_z q_\phi(z|x)log\frac{p(z)}{q_\phi(z|x)}dz \\&
    = \int_z q_\phi(z|x)logp_\theta(x|z)dz - KL[q_\phi(z|x)||p(z)]
    \\&= \mathbb{E}_{x\sim q_\phi(z|x)}[logp_\theta(x|z)] - KL[q_\phi(z|x)||p(z)]
\end{align*}
综上所述，这就是$ELBO$的数学定义：
\begin{align*}
ELBO(\theta,\sigma)= \mathbb{E}_{x\sim q_\phi(z|x)}[logp_\theta(x|z)] - KL[q_\phi(z|x)||p(z)]
\end{align*}
也可以这样表示：
\begin{align*}&
\mathcal L(\theta,\sigma)= \mathbb{E}_{x\sim q_\phi(z|x)}[logp_\theta(x|z)] - KL[q_\phi(z|x)||p(z)]\\&
\mathcal L(q) = \mathbb{E}_{q(z)}[logp(x，z)]-\mathbb{E}_{q(z)}[logq(z)]\\&
\mathcal L(q) = \mathbb{E}_{q(z)}[logp(x|z)]-KL(q(z)||p(z))
\end{align*}
\subsection{$ELBO$ 的分解形式与直观意义$^{[13]}$}
\begin{equation*}
ELBO(\theta,\sigma)= \mathbb{E}_{x\sim q_\phi(z|x)}[logp_\theta(x|z)] - KL[q_\phi(z|x)||p(z)]
\end{equation*}
\hspace*{2em} 这个形式清晰地展现了 $ELBO$ 的两个组成部分：\\
重构项（Reconstruction Term）：$\mathbb{E}_{x\sim q_\phi(z|x)}[logp_\theta(x|z)]$，衡量潜变量对数据的重构能力。\\
正则项（Regularization Term）：$KL[q_\phi(z|x)||p(z)]$，约束变分分布不要偏离先验过远。\\
对于重构项，我们希望它越大越好，如下为推导：
\begin{align*}
    Smaple\{x^1,x^2,\dots,x^m\}\, from \, P_{data}(x)
\end{align*}
\begin{align*}
    \theta^* &= arg \,\underset{\theta}{max}\prod_{i=1}^mP_\theta(x^i)= arg \,\underset{\theta}{max}\, log\prod_{i=1}^mP_\theta(x^i) \\&=
    arg \,\underset{\theta}{max}\, \prod_{i=1}^mlogP_\theta(x^i) \approx arg \,\underset{\theta}{max}\, \mathbb{E}_{x \sim P_{data}}[logP_\theta(x)] \\& = arg\, \underset{\theta}{max} \int_xP_{data}(x)logP_\theta(x)dx - \int_xP_{data}(x)logP_{data}(x)dx\\&= arg\, \underset{\theta}{max} \int_xP_{data}(x)log\frac{P_\theta(x)}{P_{data}(x)}dx \\& = arg\,\underset{\theta}{min} \,KL(P_{data}||P_\theta)
\end{align*}
即最大似然估计等价于最小化数据分布和模型分布之间的 KL 散度。\\\\
\hspace*{2em} 最大化模型生成数据的概率（似然）是目标，
重构项是生成模型 $p_\theta(x|z)$的关键部分；所以重构项越大，说明模型越会“模仿数据”，最终越接近真实分布。
这就是它与最大似然推导中 KL 散度之间的桥梁关系。

\section{$ELBO$ 的物理意义与性质}
\subsection{下界属性与推理目标}
\hspace*{2em} 在变分推断框架中，$ELBO$作为对对数边缘似然$logp_\theta(x)$的下界，有着清晰的优化目标与物理含义。根据前一章推导：
\begin{equation*}
    logp_\theta(x)= ELBO + KL[q_\phi(z|x)||p_\theta(z)]
\end{equation*}
\hspace*{2em} 由于 KL 散度非负（即 $KL[q_\phi(z|x)||p(z)] \ge 0$，我们有：
\begin{equation*}
    logp_\theta(x)\ge ELBO
\end{equation*}
\hspace*{2em} 这说明：最大化 $ELBO$ 即间接提升了模型的边缘似然 $p(x)$，从而提高模型生成数据的能力。因此，$ELBO$ 的优化可以视为一个可计算的替代目标，用于逼近原本不可解的最大似然估计问题。
\subsection{后验近似质量的衡量指标}
\hspace*{2em} $ELBO$ 不仅是一个优化目标，也可以用来衡量变分分布与真实后验之间的接近程度。由下式：
\begin{equation*}
    logp_\theta(x) - ELBO= KL[q_\phi(z|x)||p_\theta(z)]
\end{equation*}
\hspace*{2em} 我们知道，当 $q_\phi(z|x)=p_\theta(z)$时，$ELBO$ 与 $logp_\theta(x)$ 相等，此时变分推断得到最优解。因此，$ELBO$ 越接近 $logp_\theta(x)$，说明变分分布对后验的近似质量越高。
但由于$logp_\theta(x)$通常无法直接计算，这种差距本身在实际中无法精确评估，因此 $ELBO$ 也被视为一种代理目标（surrogate objective）：在不直接访问真实后验的前提下，提供一种逐步改进近似分布的方式。
\subsection{重构-正则化的权衡关系}
\subsubsection{引入超参数$\beta$$^{[14-15]}$}
\hspace*{2em} 我们可以引入通过引入超参数$\beta$显式控制 $ELBO$ 中重构项与正则项的权重，系统研究两者的权衡关系，修改标准 $ELBO$，引入系数β：
\begin{equation*}
    L(\theta,\sigma;x,\beta) = \mathbb{E}_{q_\phi(z|x)}[logp_\theta(x|z)]-\beta KL(q_\phi(z|x)||p(z))
\end{equation*}
$\beta$的作用：\\
增大 $\beta$ ⟶ 加强 KL 项的惩罚，迫使隐变量 z 更严格匹配先验 p(z)（通常为标准高斯）。\\
减小 $\beta$ ⟶ 放宽正则化约束，优先优化重构精度。\\\\
强正则化（$\beta$ > 1）：\\
\hspace*{2em} KL 项主导优化，迫使每个隐变量维度$z_i$独立且高斯分布。\\
\hspace*{2em} 效果：隐空间维度与数据中的基本视觉概念（如物体形状、颜色、位置）对齐，实现解耦表示。\\
弱正则化（$\beta$ < 1）：\\
\hspace*{2em} 重构项主导优化，允许隐变量偏离先验分布。\\
\hspace*{2em} 效果：模型更关注数据细节，生成图像更清晰，但隐变量可能耦合多个特征（可解释性下降）。\\\\
$\beta$ → 0 ⟶ 退化为自编码器（无正则化，隐空间无结构）。\\
$\beta$ → $\infty$ ⟶ 隐变量完全匹配先验（生成能力崩溃）。\\
\subsubsection{不固定$\beta$，逐渐增加 KL 的目标值$^{[16]}$}
对于原本的$\beta-ELBO$将其改写为带约束的优化问题：\\
\begin{equation*}
    \underset{\theta,\sigma}{max}\mathbb{E}_{q_\phi(z|x)}[logp_\theta(x|z)] \,\,\,\,\,s.t.\,\,\,\,\,KL(q_\phi(z|x)||p(z)) \le C
\end{equation*}
其中 $C$ 为 $KL$散度的目标容量（Capacity），随训练逐步增加（如线性增长：$C_t = C_{min}+ \frac{t}{T}(C_{max}-C_{min})$）。\\\\
通过拉格朗日乘子法，将约束优化转化为无约束问题：
\begin{equation*}
    L(\theta,\sigma;x,\beta) = \mathbb{E}_{q_\phi(z|x)}[logp_\theta(x|z)]-\lambda( KL(q_\phi(z|x)||p(z))-C)
\end{equation*}
$\lambda$ 为动态调整的拉格朗日乘子，通过梯度下降更新：
\begin{equation*}
    \lambda ← \lambda + \eta( KL(q_\phi(z|x)||p(z))-C)
\end{equation*}
其中 $\eta$ 为学习率。\\
与 $\beta-ELBO$ 的关系\\
当 $\lambda$ 收敛时，等价于固定 $\beta = \lambda$，但 $KL$ Capacity 的 $\lambda$ 是自适应调整的，且 $C$ 逐步增加，避免了静态 $\beta$ 的局限性。\\\\
Capacity 调度策略:\\
线性增长：$C_t = C_{min}+ \frac{t}{T}(C_{max}-C_{min})$\\
余弦退火：$C_t = C_{min}+ \frac{t}{T}(C_{max}-C_{min})(1-cos(\frac{t\pi}{T}))$
目标值选择：\\
$C_{min}$：初始容量（如 0），允许模型早期自由学习。\\
$C_{max}$：最终容量（如 20$\sim$50），控制解耦强度。\\
\hspace*{2em} 训练初期，$KL$ 容量 $C$ 较小，模型优先优化重构项，学习有意义的隐变量表示。
训练后期，逐步增大 $C$，迫使隐变量分布逐渐逼近先验，实现解耦。
在训练期间会进行自适应平衡：拉格朗日乘子 $\lambda$ 自动调节 $KL$ 项的权重，避免人工调参，解耦与生成质量更优。

\subsection{多任务视角与平衡策略$^{[17]}$}
\hspace*{2em} 通过显式解耦总相关性（Total Correlation, $T$C） 来增强隐变量的独立性。结合 Capacity Constraint（动态调整 $KL$ 约束），可以更灵活地平衡解耦与生成质量。\\
总相关性（$TC$）的定义:\\
$TC$ 衡量隐变量维度之间的依赖性：
\begin{equation*}
    TC(q_\phi(z)) = KL(q_\phi(z)||\prod_{j=1}^d q_\phi(z_j))
\end{equation*}
$q_\phi(z)$是隐变量的聚合后验\\
$\prod_{j=1}^d q_\phi(z_j)$是假设各维度独立时的联合分布\\
$TC = 0$ 表示所有维度完全独立\\\\
显式优化 $TC$ 项：
\begin{equation*}
    L_{TC} = \mathbb{E}_{q_\phi(z|x)}[logp_\theta(x|z)] - \alpha KL(q_\phi(z|x)||p(z))-\beta TC(q_\phi(z))
\end{equation*}
$\alpha$ 控制先验匹配强度（默认 $\alpha$=1）。\\
$\beta$ 控制解耦强度（通常 $\beta$>1）。\\\\
\hspace*{2em} 直接计算 $TC$ 需要估计$q_\phi(z)$ ，但其解析形式未知。TC-VAE采用两种估计策略：\\
(1)密度比估计:\\
利用判别器 D(z) 区分样本来源:
\begin{equation*}
    TC(q_\phi(z)) \approx \mathbb{E}_{q_\phi(z)}[log\frac{D(z)}{1-D(z)}]
\end{equation*}
正样本：从$q_\phi(z)$采样（即$z \sim q_\phi(z), x\sim p_{data}(x)$）。\\
负样本：从$\prod_{j=1}^dq_\phi(z_j)$采样（各维度独立重排）。\\\\
(2) Minibatch 加权采样:\\
通过 minibatch 近似$q_\phi(z)$：
\begin{equation*}
    TC \approx \mathbb{E}_{q_\phi(z)}[logq_\phi(z)] - \sum_{j=1}^d\mathbb{E}_{q_\phi(z)}[logq_\phi(z_j)]
\end{equation*}
计算复杂度高，但对小规模数据可行。\\\\
动态约束的数学形式\\
将 TC 和 KL 项分别施加容量约束：\\
\begin{equation*}
    L_{TC+Capacity} = \mathbb{E}_{q_\phi(z|x)}[logq_\phi(x|z)] - \lambda_1(KL(q_\phi(z|x)||p(z))-C_{KL}) - \lambda_2(TC(q_\phi(z))-C_{TC})
\end{equation*}
$C_{KL}$和$C_{TC}$是逐步增加的容量目标。\\$\lambda_1,\lambda_2$为拉格朗日乘子，通过梯度上升更新：
\begin{equation*}
    \lambda_i←\lambda_i+\eta(Constraint_i-C_i)
\end{equation*}
容量调度策略:\\
$KL$ 容量（$C_{KL}$）：从 0 线性增长到目标值（如 20）。\\
$TC$ 容量（$C_{TC}$）：从 0 增长到更低的值（如 5），因为 $TC$ 的解耦需求通常弱于全局 $KL$。\\\\
\hspace*{2em} 其可以通过 TC 项直接优化隐变量独立性。并且可以动态平衡容量约束，实现训练过程的自动调参。
\section{$ELBO$ 的优化方法}
\hspace*{2em} 变分推断的核心任务是最大化 $ELBO$，从而近似计算后验分布。由于 $ELBO$ 涉及期望运算与 KL 散度，其优化在理论与工程实践中均面临挑战。本章将综述常见的 $ELBO$ 优化方法，涵盖解析推导、数值估计与现代深度学习中的变分推断技术。
\subsection{坐标上升法：Mean-Field 变分推断$^{[18-20]}$}
Mean-Field 变分推断（Mean-Field Variational Inference, MFVI）是变分推断中最经典的方法之一，其核心思想是通过坐标上升法（Coordinate Ascent）迭代优化 $ELBO$。
假设隐变量$\mathbf{z} = (z_1,z_2,\dots,z_K)$的各维度完全独立，变分分布分解为：
\begin{equation*}
   q(\mathbf{z}) = \prod_{k=1}^K q_k(z_k)
\end{equation*}
每个 $q_k(z_k)$ 是边缘分布（如高斯、伯努利等）。\\
该假设简化了计算，但牺牲了隐变量间的相关性建模能力。\\
在 Mean-Field 假设下，$ELBO$ 可写为：
\begin{align*}
   \mathcal L(q) &= \mathbb{E}_{q(\mathbf{z})}[\log p(\mathbf{x},\mathbf{z})] - \mathbb{E}_{q(\mathbf{z})}[\log q(\mathbf{z})]\\&=\mathbb{E}_{q(\mathbf{z})}[\log p(\mathbf{x},\mathbf{z})] - \sum_{k=1}^K \mathbb{E}_{q_k(z_k)}[\log q_k(z_k)]
\end{align*}
\hspace*{2em} 固定其他维度，优化单维度$q_k$
 对每个$q_k(z_k)$，固定其他$q_{-k}(z_{-k}) = \prod_{j\neq k}q_j(z_j)$，$ELBO$ 中仅与 $q_k$相关的项为
\begin{align*}
\mathcal L(q_k) &= \mathbb{E}_{q_k(z_k)}\mathbb{E}_{q_{-k}(\mathbf{z}_{-k})}[\log p(\mathbf{x},\mathbf{z})] - \mathbb{E}_{q_k(z_k)}[logq_k(z_k)] + const \\ &=\mathbb{E}_{q_k(z_k)}\mathbb{E}_{q_{-k}(\mathbf{z}_{-k})}[\log p(\mathbf{x},\mathbf{z})] - log\tilde{p}_k(z_k)
\\&= - KL(q_k(z_k)||\tilde{p}_k(z_k))
\end{align*}
最大化 $ELBO$ 等价于最小化 KL 散度，因此最优解为：
\begin{align*}
q_k^*(z_k) \propto \exp\left(\mathbb{E}_{q_{-k}(\mathbf{z}_{-k})}[\log p(\mathbf{x},\mathbf{z})]\right)
\end{align*}
对每个维度 $k$，按以下步骤更新：\\
计算期望：
\begin{equation*}
\log q_k(z_k) ← \mathbb{E}_{q_{-k}(\mathbf{z}_{-k})}[\log p(\mathbf{x},\mathbf{z})] + \text{const}
\end{equation*}
归一化：
\begin{equation*}
\log q_k(z_k) = \frac{exp(\mathbb{E}_{q_{-k}}[logp(\mathbf{x},\mathbf{z})])}{\int exp(\mathbb{E}_{q_{-k}}[logp(\mathbf{x},\mathbf{z})])dz_k}
\end{equation*}
\hspace*{2em} Mean-Field 变分推断中的坐标上升法通过迭代优化单变量边缘分布最大化 $ELBO$, 这种方法适用于具有 共轭先验 的模型，如变分贝叶斯 LDA、GMM、贝叶斯回归等，更新过程高效且收敛稳定。
\subsection{随机梯度变分推断（SGVI）：Black-Box 方法$^{[21-22]}$}
\hspace*{2em} 针对复杂模型、非共轭结构以及神经网络等不具备解析解的情况，采用 黑盒变分推断（Black-Box Variational Inference），即通过采样与梯度估计进行优化。\\
梯度估计公式：
$ELBO$ 关于变分参数$\phi$的梯度可写为：
\begin{equation*}
\nabla_\phi \mathcal L(\phi) = \nabla_\phi \mathbb{E}_{q_\phi(z)}[log(x,z)-logq_\phi (z)]
\end{equation*}
由于期望不可解析，可通过蒙特卡洛采样估计：
\begin{equation*}
\nabla_\phi \mathcal L \approx \frac{1}{K}\sum_{k=1}^K  \nabla_\phi \,logq_\phi (z)\, [logp(x,z)-logq_\phi (z)]
\end{equation*}
这种方法通用性强，但方差较高，常需方差控制技术如 control variates 辅助。
\subsection{重参数化技巧（Reparameterization Trick）$^{[23]}$}
\hspace*{2em} 引入了重参数化技巧，可以极大降低梯度估计的方差，使得 $ELBO$ 在神经网络中可直接通过反向传播进行优化。\\
将$z\sim q_\phi(z,x)$转换为一个确定函数：
\begin{equation*}
z = g_\phi(\epsilon),\epsilon\sim p(\epsilon)
\end{equation*}
例如对高斯分布：
\begin{equation*}
z = \mu(x)+\sigma(x)\epsilon,\epsilon\sim \mathcal N(0,1)
\end{equation*}
此时 $ELBO$ 的梯度变为：
\begin{equation*}
\nabla_\phi \mathcal L \approx \frac{1}{K}\sum_{k=1}^K  \nabla_\phi \,logp(x,g_\phi(\epsilon_k))-logq_\phi(g_\phi(\epsilon_k)|x)
\end{equation*}
该技巧目前被广泛用于 VAE、Bayesian Neural Network 等深度模型中。
\subsection{自然梯度变分推断（Natural Gradient VI）$^{[24-25]}$}
\hspace*{2em} 标准梯度在高维概率分布参数空间中更新效率低下。自然梯度（Natural Gradient） 在信息几何上考虑变分分布之间的真实“距离”，提高了收敛速度和稳定性。\\
核心思想：\\
使用变分参数的 Fisher 信息矩阵 $F(\phi)$，更新方向变为：
\begin{equation*}
\tilde\nabla_\phi \mathcal L = F(\phi)^{-1}\nabla_\phi \mathcal L
\end{equation*}
\hspace*{2em} 在指数族分布中，自然梯度往往有简洁的闭式形式；在深度学习中，也有 K-FAC 等近似方法用于高效计算。
\section{$ELBO$ 的典型应用场景}
\hspace*{2em} $ELBO$（Evidence Lower Bound）不仅是变分推断的核心目标函数，更是多个现代概率模型和生成模型中的理论基石。由于其能有效逼近后验分布与边缘似然，$ELBO$ 被广泛用于各种贝叶斯模型、深度生成模型、聚类模型等领域。以下将列举 $ELBO$ 的典型应用场景，并简要介绍其在具体模型中的作用。
\subsection{变分自编码器（Variational Autoencoders, VAE）$^{[26-28]}$}
\hspace*{2em} VAE 是 $ELBO$ 在深度学习中最具代表性的应用。其基本思想是引入隐变量 $z$ 来建模数据 $x$ 的生成过程，通过最大化 $ELBO$ 学习生成模型的参数。
\subsubsection{VAE 的结构组成}
编码器（Encoder）:\\
\hspace*{2em} 将输入 $x$ 映射为潜在变量 $z$ 的分布：
\begin{equation*}
q_\phi(z|x)
\end{equation*}
常假设这个分布是正态分布 $\mathcal N(\mu,\sigma^2)$由神经网络输出$\mu(x),\sigma(x)$。 \\
\hspace*{2em}什么是映射呢？我们不再像传统自编码器那样把输入压缩成一个固定的向量$z$，而是将 $x$ 映射成一个概率分布（比如正态分布），然后从这个分布中采样得到 $z$。\\
传统 AutoEncoder 的做法：\\
\hspace*{2em} 给定一个图片 $x$，编码器输出一个向量 $z = f(x)$。这个 $z$ 是固定的——每次输入 $x$，输出都是一样的 $z$。\\
VAE 的做法：\\
\hspace*{2em} 给定一个图片 $x$，编码器输出的是两个向量：$\mu(x),\sigma(x)$,它们定义了一个分布$q_\phi(z|x) = \mathcal N(z;\mu,\sigma^2)$我们从这个分布中随机采样一个 $z$，再输入给解码器重建。这个“随机性”可以让模型更有泛化能力，也让它能学习生成模型。\\
\hspace*{2em} 编码器不是输出一个确定的潜在向量$z$，而是输出一个高斯分布（用均值和标准差描述），表示“这张图片在潜在空间中可能落在什么位置”，从而可以采样多个$z$，用于生成或重建。\\\\
解码器（Decoder）:\\
根据潜在变量 $z$ 重建原始输入：
\begin{equation*}
p_\theta(x|z)
\end{equation*}
\hspace*{2em} 这是一个从潜在空间回到数据空间的映射，也是一个神经网络。\\
\hspace*{2em} 解码器是一个神经网络，学习的是$p_\theta(x|z)$也就是在给定潜在变量$z$的条件下，输入数据$x$的概率。就是希望这个解码器能够生成尽量接近原始输入$x$的输出。\\
传统 AutoEncoder 的做法：\\
\hspace*{2em} 将压缩后的向量还原为原始数据\\
VAE 的做法：\\
\hspace*{2em} 将从分布中采样的$z$还原为原始数据，并保证可以“生成”新样本。它还可以通过在两个$z_1,z_2$之间线性插值，然后解码，得到连续变化的图像（比如从“3”慢慢变成“8”），并在潜在空间中轻微扰动$z$，可得到风格变化图像。\\
\hspace*{2em} 所以VAE可以输入一个样本，输出一个和样本相似的东西。
\subsubsection{为什么叫“变分”自编码器}
\hspace*{2em} VAE 使用了变分推断思想，通过优化 $ELBO$ 目标函数，来近似最大化数据的对数似然$logp(x)$\\
什么是$p(x)$,我们为什么需要这个$p(x)$?\\
\hspace*{2em} $p(x)$是模型生成某个具体样本 $x$的“概率”。它告诉我们，模型认为这个样本有多“合理”或“典型”。\\
为什么要关心$p(x)$?\\
\hspace*{2em} 我们训练生成模型的目标，实际上是让模型能够给训练数据（和相似数据）更高的概率，即模型能更好地理解数据的分布。这意味着，最大化所有训练样本的 $p(x)$，使模型“学会”训练集的分布。模型能用它学到的概率规则，生成类似训练数据的样本。
$p(x)$通常很小，怎么办？\\
\hspace*{2em} 数据空间（比如图像空间）极其庞大且稀疏，一个具体的样本$x$的概率往往非常小，甚至趋近于0。但这没关系，因为我们比较的是不同样本的概率大小关系，而不是概率绝对值。所以通常用 对数概率来表示，更易于比较和优化。\\\\
\hspace*{2em} 在网络中$p(x)$不只是一个评价网络好坏的标准，我们还使用它训练我们的网络，但
训练时不直接用$p(x)$计算梯度，但它太难算了，所以我们训练用的是$ELBO$。$ELBO$ 是一个可计算、可导的目标函数，方便我们用梯度下降优化。在 VAE 训练里，我们通过计算 $ELBO$ 的梯度来更新编码器和解码器的参数。
\subsection{贝叶斯神经网络（Bayesian Neural Networks, BNN）$^{[29-30]}$}
\hspace*{2em}在普通的神经网络中，训练的目标是找到一组最优参数 $\theta$（例如权重、偏置），使得网络在训练集上表现最好。但是它没有考虑一个问题，如果训练数据不够充分、带有噪声，为什么我们要那么“确信”这组参数是最优的？\\
\hspace*{2em}所以贝叶斯神经网络BNN 不是寻找“最优权重”，而是为每个权重分配一个概率分布，AVE是将输入概率化的话，那么BNN就是将权重概率化。\\
训练时，我们的思想和VAE一样，使用变分推断。\\
假设我们有训练数据集：
\begin{equation*}
D=\{(x^{(i)},y^{(i)})\}_{i=1}^N
\end{equation*}
我们关心的是，给定数据$D$，推断参数$\theta$的后验分布：
\begin{equation*}
p(\theta|D) = \frac{p(D|\theta)}{p(D)}
\end{equation*}
\hspace*{2em}前面也说了这个后验通常非常难直接计算（高维积分、没有闭式解），所以我们采取变分推断方法进行近似。\\
我们构造一个可学习的近似分布$q_\phi(\theta)$，其目标是：
\begin{align*}
\log p(D)
&= \log \int p(D \mid w) p(w) \, dw \\
&= \log \int q(w) \frac{p(D \mid w) p(w)}{q(w)} \, dw \\
&\geq \int q(w) \log \frac{p(D \mid w) p(w)}{q(w)} \, dw \\
&= \mathbb{E}_{q(w)}[\log p(D \mid w)] - \mathrm{KL}(q(w) \| p(w)) \\
&= \mathcal{L}(q)
\end{align*}
$p(D|w)$：似然项（神经网络的预测概率）\\
$p(w)$：先验分布（常为标准正态）\\
$q(w)$：近似后验分布（可参数化为高斯分布）\\
$KL(q(w)||p(w))$正则项，鼓励 $q(w)$ 不偏离先验\\

这意味着我们希望 $q_\phi(\theta)$逼近真实后验。由于$p(\theta|D)$包含不可计算的$p(D)$，这个目标等价于最大化$ELBO$,也就是 BNN 的训练目标:
\begin{equation*}
logp(D) \ge \mathbb{E}_{q_\phi(\theta)}[logp(D|\theta)]-KL(q_\phi(\theta)||p(\theta))
\end{equation*}
\hspace*{2em}在预测的时候我们就不要算复杂的积分了，只需要从$q_\phi(\theta)$中采样多组$\theta$，对每组$\theta$计算$p(y^*,x^*,\theta)$最终取平均，获得预测分布。\\
预测时综合所有可能的模型（权重），自然具备“模型集成”和“不确定性量化”的能力。
\subsection{潜在主题模型（如 Latent Dirichlet Allocation, LDA）$^{[31]}$}
\hspace*{2em}LDA 是文本分析中的经典模型，用于从文档中挖掘潜在主题。其模型包含文档-主题分布和主题-词分布等多个隐变量，直接推断后验分布不可行。
使用我们必须使用变分方法了，可以使用 Mean-field 结构近似各隐变量后验，然后最大化 $ELBO$ 以优化主题分布与词分布。\\
\hspace*{2em}变分 LDA 是早期变分推断的重要成果，也是自然语言处理领域中 $ELBO$ 的成功应用实例。
\subsection{贝叶斯高斯混合模型（Bayesian GMM）$^{[32]}$}
\hspace*{2em}高斯混合模型是无监督学习中的经典聚类模型，引入贝叶斯推断后，对每个样本的类别分布和模型参数引入不确定性。\\
$ELBO$ 表达式：
\begin{align*}
\log p(X)
&= \log \int p(X, Z, \pi, \mu, \Lambda) \, dZ \, d\pi \, d\mu \, d\Lambda \\
&\geq \mathbb{E}_{q(Z, \pi, \mu, \Lambda)} \left[ \log \frac{p(X, Z, \pi, \mu, \Lambda)}{q(Z, \pi, \mu, \Lambda)} \right] \\
&= \mathcal{L}(q)
\end{align*}
\hspace*{2em}其中，$z$ 表示隐类别标签，$\pi, \mu, \Lambda$为混合参数。
$ELBO$ 优化过程可由变分 EM 或随机推断方法实现，使模型在聚类任务中更具鲁棒性与泛化能力。

\subsection{强化学习中的贝叶斯策略优化$^{[33-34]}$}
\hspace*{2em}在强化学习中，策略的不确定性建模成为提升探索效率的重要手段。基于 $ELBO$ 的贝叶斯策略优化方法，如 Variational Policy Search，使用变分方法估计策略后验。\\
$ELBO$ 表达式：
\begin{align*}
\mathcal{L}(q)
&= \mathbb{E}_{q(Z)}[\log p(X | Z, \mu, \Lambda)]
+ \mathbb{E}_{q(Z)q(\pi)}[\log p(Z | \pi)]
+ \mathbb{E}_{q(\pi)}[\log p(\pi)] \\
&\quad + \mathbb{E}_{q(\mu, \Lambda)}[\log p(\mu, \Lambda)]
- \mathbb{E}_{q(Z)}[\log q(Z)]
- \mathbb{E}_{q(\pi)}[\log q(\pi)]
- \mathbb{E}_{q(\mu, \Lambda)}[\log q(\mu, \Lambda)]
\end{align*}
这种方法可提升策略稳定性，并实现置信度感知的行为决策。\\\\\\


\hspace*{2em}$ELBO$ 已成为现代概率建模的统一语言与优化目标。无论是在生成建模（VAE）、后验估计（BNN）、文本聚类（LDA）、还是策略学习（RL）中，$ELBO$ 提供了一种稳定、理论完备的优化框架。随着模型复杂性的增加，对 $ELBO$ 更高效、更精确的估计与优化方法也成为研究前沿。

\section{$ELBO$ 的改进与前沿方向}
\hspace*{2em}尽管 $ELBO$ 在变分推断中发挥着核心作用，但其存在一些固有局限例如下界较松、优化过程易陷入局部最优等。因此，近年来针对 $ELBO$ 的精度与适应性，提出了多种改进策略与扩展方法，以提升近似后验的质量和模型表达能力。
\subsection{更紧的下界$^{[35]}$}
\hspace*{2em}传统 $ELBO$ 提供的是边缘似然$logp(x)$的一个下界，但该下界可能偏差较大。IWAE 提出使用多个样本进行重要性加权，得到更紧的下界：
\begin{equation*}
\mathcal{L}_K = \mathbb{E}_{z_1,\dots,z_K\sim q(z|x)}[log(\frac{1}{K})\sum_{k=1}^K\frac{p(x,z_k)}{q(z_k|x)}]
\end{equation*}
\hspace*{2em}当$K=1$时，该表达式退化为传统 $ELBO$；当$K$增大时，$\mathcal{L}_K$更接近$logp(x)$，从而提供更精确的优化目标。IWAE 改进了生成模型的训练稳定性与生成样本的多样性。
\subsection{替代 KL 散度的变分目标}
\hspace*{2em}$ELBO$ 的推导依赖于KL散度尤其是(尤其是$KL(q(z)||p(z|x))$)，但KL散度具有“零覆盖问题”，即若
$q(z)$ 不覆盖$p(z|x) $的支持区域，其值将趋于无穷，导致训练不稳定。为此，可以引入了更一般的散度度量。\\
\subsubsection{$\chi$散度$^{[36]}$}
$\chi$散度是 $f$-散度的一种形式，其定义为：
\begin{equation*}
\chi^2(q||p) = \int(\frac{q(z)}{p(z)}-1)^2p(z)dz
\end{equation*}
相比 KL 散度，$\chi$散度对$p(z)$的尾部行为更加敏感，优化过程中鼓励更宽泛的支持区域。
\subsubsection{R-$ELBO$$^{[37]}$}
基于 Rényi 散度的目标函数：
\begin{equation*}
\mathcal{L}_\alpha = \frac{1}{1-\alpha}log\mathbb{E}_{q(z)}[(\frac{p(x,z)}{q(z)})^{1-\alpha}]
\end{equation*}
其中$\alpha$是调节参数，当$\alpha→1$时，$\mathcal{L}_\alpha$收敛为标准 $ELBO$。调节$\alpha$可控制偏差与方差的权衡。\\\\\\
\hspace*{2em}$ELBO$ 虽为经典目标，但仍不断演化。未来，如何在表达能力、计算效率与收敛性之间平衡，将是 $ELBO$ 研究的重要方向。


\section{总结}
\hspace*{2em}本文系统回顾了 $ELBO$（Evidence Lower Bound, 证据下界）作为变分推断核心目标函数的理论基础与实际应用价值。从基本数学推导出发，我们展示了 $ELBO$ 如何通过 Jensen 不等式和 KL 散度的重写形式，构建出一种可优化的、变分近似后验的评估指标。通过最大化 $ELBO$，我们既可以近似最大化观测数据的对数边际似然，又能最小化后验分布与近似分布之间的差异。\\

$ELBO$ 在变分推断、生成模型、贝叶斯神经网络、聚类等众多任务中展现出广泛应用价值，其优化技术也经历了从解析式更新、重参数化梯度、自然梯度到黑盒变分推断等多种演化形式。此外，围绕其表达能力、优化稳定性与近似精度的不足，研究者提出了诸如 IWAE、$\chi$-VI 等改进方向，推动了高维复杂模型的可扩展推断。







\begin{thebibliography}{00}
\bibitem{jordan1999vi}
M.~I. Jordan, Z. Ghahramani, T. Jaakkola, and L.~K. Saul,
An introduction to variational methods for graphical models, in
{\it Machine Learning}, vol.~37, no.~2, pp. 183--233, 1999.

\bibitem{zhang2017advances}
C. Zhang, J. Butepage, H. Kjellstrom, and S. Mandt,
Advances in variational inference,
{\it IEEE Transactions on Pattern Analysis and Machine Intelligence}, 2017.

\bibitem{kipf2016}
T. N. Kipf and M. Welling,
Variational Graph Auto-Encoders,
{\it arXiv preprint arXiv:1611.07308}, 2016.

\bibitem{kim2018}
Y. Kim, S. Wiseman, and A. C. Miller,
Semi-Amortized Variational Autoencoders,
{\it Proc. Int. Conf. on Machine Learning (ICML)}, Stockholm, Sweden, 2018.

\bibitem{dempster1977em}
A.~P. Dempster, N.~M. Laird, and D.~B. Rubin,
Maximum likelihood from incomplete data via the EM algorithm,
{\it Journal of the Royal Statistical Society, Series B}, vol.~39, no.~1, pp. 1–22, 1977.

\bibitem{hunlearn2019em}
Hun-Learning,
EM Algorithm for Latent Variable Models – Derivation of the ELBO,
{\it Hun-Learning Technical Blog}, 2019.

\bibitem{daza2020variational}
D. Daza,
Expectation Maximization for Latent Variable Models – A Variational Perspective,
{\it Daza’s Machine Learning Notes}, 2020.

\bibitem{blei2017vi}
D.~M. Blei, A. Kucukelbir, and J.~D. McAuliffe,
Variational inference: A review for statisticians,
{\it Journal of the American Statistical Association}, vol.~112, no.~518, pp. 859–877, 2017.

\bibitem{bishop2006prml}
C.~M. Bishop and N.~M. Nasrabadi,
Variational inference and the evidence lower bound, in
{\it Pattern Recognition and Machine Learning}, Springer, 2006, ch.~10.

\bibitem{odaibo2019vae}
S. Odaibo,
Tutorial: Deriving the Standard Variational Autoencoder (VAE) Loss Function,
{\it arXiv preprint}, July 2019.

\bibitem{bishop2006prml}
C.~M. Bishop and N.~M. Nasrabadi,
Variational inference and the evidence lower bound, in
{\it Pattern Recognition and Machine Learning}, Springer, 2006, ch.~10.

\bibitem{blei2017vi}
D.~M. Blei, A. Kucukelbir, and J.~D. McAuliffe,
Variational inference: A review for statisticians,
{\it Journal of the American Statistical Association}, vol.~112, no.~518, pp. 859–877, 2017.

\bibitem{bishop2006}
C. M. Bishop,
Pattern Recognition and Machine Learning,
Springer, 2006, Chapter 10.

\bibitem{higgins2017}
I. Higgins, L. Matthey, A. Pal, C. Burgess, X. Glorot, M. Botvinick, S. Mohamed, and A. Lerchner,
β-VAE: Learning Basic Visual Concepts with a Constrained Variational Framework,
{\it Proc. Int. Conf. on Learning Representations (ICLR)}, Toulon, France, 2017.

\bibitem{burgess2018}
C. P. Burgess, I. Higgins, A. Pal, L. Matthey, N. Watters, G. Desjardins, and A. Lerchner,
Understanding disentangling in β-VAE,
{\it arXiv preprint arXiv:1804.03599}, 2018.

\bibitem{rezende2018}
D. J. Rezende and S. Mohamed,
Tighter variational bounds are not necessarily better,
{\it Proc. Int. Conf. on Machine Learning (ICML)}, Stockholm, Sweden, 2018.

\bibitem{lyu2021}
D. Lyu, Z. Ghahramani, and F. Hernández-Lobato,
Learning Interpretable Deep Disentangled Representations with Hyperbolic VAE,
{\it Proc. Adv. Neural Inf. Process. Syst. (NeurIPS)}, Virtual Conference, 2021.

\bibitem{parisi1988}
G. Parisi,
Statistical Field Theory,
Addison-Wesley, 1988, Chapter 5.

\bibitem{jordan1999}
M. I. Jordan, Z. Ghahramani, T. S. Jaakkola, and L. K. Saul,
An introduction to variational methods for graphical models,
{\it Machine Learning}, vol. 37, no. 2, pp. 183–233, 1999.

\bibitem{bishop2006}
C. M. Bishop,
Pattern Recognition and Machine Learning,
Springer, 2006, Chapter 10.

\bibitem{ranganath2014}
R. Ranganath, S. Gerrish, and D. M. Blei,
Black Box Variational Inference,
{\it Proc. 17th Int. Conf. on Artificial Intelligence and Statistics (AISTATS)}, Reykjavik, Iceland, 2014.

\bibitem{hoffman2013}
M. D. Hoffman, D. M. Blei, C. Wang, and J. Paisley,
Stochastic Variational Inference,
{\it Journal of Machine Learning Research}, vol. 14, pp. 1303–1347, 2013.

\bibitem{kingma2014}
D. P. Kingma and M. Welling,
Auto-Encoding Variational Bayes,
{\it Proc. Int. Conf. on Learning Representations (ICLR)}, Banff, Canada, 2014.

\bibitem{mnih2014}
A. Mnih and D. J. Rezende,
Variational Inference for Monte Carlo Objectives,
{\it Proc. Int. Conf. on Machine Learning (ICML)}, Beijing, China, 2016.

\bibitem{kucukelbir2017}
A. Kucukelbir, D. Tran, R. Ranganath, A. Gelman, and D. M. Blei,
Automatic Differentiation Variational Inference,
{\it Journal of Machine Learning Research}, vol. 18, no. 14, pp. 1–45, 2017.

\bibitem{kingma2014auto}
D.~P. Kingma and M. Welling,
Auto‑Encoding Variational Bayes, in
{\it Proc. 2nd Int. Conf. on Learning Representations (ICLR)}, 2014.

\bibitem{kingma2019intro}
D.~P. Kingma and M. Welling,
An introduction to variational autoencoders,
{\it Foundations and Trends in Machine Learning}, vol.~12, no.~4, pp. 307--392, 2019.

\bibitem{rezende2014}
D. J. Rezende, S. Mohamed, and D. Wierstra,
Stochastic Backpropagation and Approximate Inference in Deep Generative Models,
{\it Proc. Int. Conf. on Machine Learning (ICML)}, Beijing, China, 2014.

\bibitem{mackay1992}
D. J. C. MacKay,
A Practical Bayesian Framework for Backpropagation Networks,
{\it Neural Computation}, vol. 4, no. 3, pp. 448–472, 1992.

\bibitem{hinton1993}
G. E. Hinton and D. van Camp,
Keeping the Neural Networks Simple by Minimizing the Description Length of the Weights,
{\it Proc. 6th Annu. Conf. on Computational Learning Theory (COLT)}, Santa Cruz, USA, 1993.

\bibitem{blei2003}
D. M. Blei, A. Y. Ng, and M. I. Jordan,
Latent Dirichlet Allocation,
{\it Journal of Machine Learning Research}, vol. 3, pp. 993–1022, 2003.

\bibitem{rasmussen2000}
C. E. Rasmussen,
The Infinite Gaussian Mixture Model,
{\it Proc. Adv. Neural Inf. Process. Syst. (NeurIPS)}, Denver, USA, 2000.

\bibitem{peters2010}
J. Peters, K. Mülling, and Y. Altun,
Relative Entropy Policy Search,
{\it Proc. AAAI Conf. on Artificial Intelligence}, Atlanta, USA, 2010.

\bibitem{neumann2009}
G. Neumann and J. Peters,
Fitted Q-iteration by Advantage Weighted Regression,
{\it Proc. Adv. Neural Inf. Process. Syst. (NeurIPS)}, Vancouver, Canada, 2009.

\bibitem{burda2016}
Y. Burda, R. Grosse, and R. Salakhutdinov,
Importance Weighted Autoencoders,
{\it Proc. Int. Conf. on Machine Learning (ICML)}, New York, USA, 2016.

\bibitem{pearson1900}
K. Pearson,
On the Criterion that a Given System of Deviations from the Probable in the Case of a Correlated System of Variables is Such that It Can Be Reasonably Supposed to Have Arisen from Random Sampling,
{\it Philosophical Magazine}, vol. 50, no. 302, pp. 157–175, 1900.


\bibitem{renyi1961}
A. Rényi,
On Measures of Entropy and Information,
{\it Proc. 4th Berkeley Symp. on Mathematical Statistics and Probability}, vol. 1, pp. 547–561, 1961.

\end{thebibliography}

\end{document}
